\documentclass[11pt,a4paper]{ltjsarticle}

\usepackage{luatexja-fontspec}
\usepackage[margin=25mm]{geometry}
\usepackage{amsmath,amssymb,amsthm,mathtools}
\usepackage{booktabs}
\usepackage{tabularx}
\usepackage{tikz-cd}
\usepackage{enumitem}
\usepackage{hyperref}

\hypersetup{
    colorlinks=true,
    linkcolor=blue,
    citecolor=blue,
    urlcolor=blue,
    pdftitle={構造変換の圏論的基礎：意味の場としてのブリッジ、アスペクト的意味の場による境界条件、およびグリーン関数との同型性},
    pdfauthor={太田 寛 (Hiroshi Ota)}
}

% 定理環境（英語版 Paper 3 と完全対応）
\theoremstyle{definition}
\newtheorem{definition}{定義}[section]
\newtheorem{axiom}{公理}[section]

\theoremstyle{plain}
\newtheorem{theorem}{定理}[section]
\newtheorem{proposition}{命題}[section]
\newtheorem{lemma}{補題}[section]
\newtheorem{corollary}{系}[section]

\theoremstyle{remark}
\newtheorem{remark}{注記}[section]

% 証明環境名の日本語化
\renewcommand{\proofname}{\textbf{証明の概略}}

% カスタムマクロ
\newcommand{\Ob}{\mathrm{Ob}}
\newcommand{\Valid}{\mathrm{Valid}}
\newcommand{\id}{\mathrm{id}}
\newcommand{\colim}{\operatorname{colim}}
\newcommand{\TwoSet}{\mathbf{2\text{-}Set}}
\newcommand{\Set}{\mathbf{Set}}
\newcommand{\CFORM}{\mathcal{C}_{\mathrm{FORM}}}
\newcommand{\CSUB}{\mathcal{C}_{\mathrm{SUB}}}
\newcommand{\CBridge}{\mathcal{C}_{\mathrm{Bridge}}}
\newcommand{\CAspect}{\mathcal{C}_{\mathrm{Aspect}}}
\newcommand{\CSrc}{\mathcal{C}_{\mathrm{Src}}}
\newcommand{\CTarg}{\mathcal{C}_{\mathrm{Targ}}}
\newcommand{\MSrc}{\mathcal{M}_{\mathrm{Src}}}
\newcommand{\MTarg}{\mathcal{M}_{\mathrm{Targ}}}
\newcommand{\MAspect}{\mathcal{M}_{\mathrm{Aspect}}}
\newcommand{\MSub}{\mathcal{M}_{\mathrm{Sub}}}
\newcommand{\BTrans}{\mathcal{B}_{\mathrm{Trans}}}
\newcommand{\BAspect}{\mathcal{B}_{\mathrm{Aspect}}}
\newcommand{\Eexec}{\mathcal{E}_{\mathrm{exec}}}

\title{\textbf{構造変換の圏論的基礎：\\
意味の場としてのブリッジ、アスペクト的意味の場による境界条件、およびグリーン関数との同型性}}

\author{太田 寛（Hiroshi Ota）\thanks{10.5281/zenodo.24431232: \url{https://doi.org/10.5281/zenodo.24431232} \\ Homepage: \url{https://www.kae-made.jp}, Research Notes: \url{https://note.com/kae_made/membership}} \\
\small Independent Researcher / Knowledge \& Experience\\
\small \texttt{master@kae-made.jp}}

\date{\today}

\begin{document}

\maketitle

\begin{abstract}
前三稿（Paper 0~\cite{ota2026paper0}, Paper 1~\cite{ota2026paper1}, Paper 2~\cite{ota2026paper2}）において、著者は関係を実在の基底とする存在論と複数の「意味の場（Fields of Sense）」からなるユニバースを定式化し、単一の意味の場における静的な意味構造（概念情報モデル：商 2-圏 $\mathcal{C} \equiv CG/\mathcal{R}$）ならびに動的振る舞い（状態モデル、非循環データフロー型アクションモデルによる対称モノイド圏 $\mathcal{DF}$、および形式への言及 $\CFORM$ に基づく普遍的実行セマンティクス $\Eexec$）が、物理学の理論記述と完全に同型な閉じた数理体系をなすことを証明した。しかし、現実世界の現象や知的・工学的営為は孤立した単一の意味の場のみで完結するものではなく、並立する複数の意味の場同士の相互作用、ならびに「ある意味の場のスキーマ（形式）を、別の意味の場のスキーマ制約に従う具体的な状態（インスタンスモデル）として構成する次元転換」を伴って立ち現れる。

本稿（Paper 3）は、視座を単一の意味の場から「複数の意味の場間の対応付けと境界値変換」へと拡張し、次の4つの根源的命題を証明することを目的とする。第一に、意味の場同士を結ぶ\textbf{ブリッジ（Bridge）は特権的な外部装置ではなく、それ自体が固有の概念情報モデル・状態モデル・アクションモデル・実行セマンティクスを備えた独立した「意味の場」として概念モデリング体系そのもので記述可能である}こと（自己閉包性と自己相似性）を証明し、対向ドメインのクラスを映し出す\textbf{プロジェクションクラス}と意味的フレーズを持たない純粋な構造対応射としての\textbf{ファンクター}を定式化する。第二に、概念モデリングの「形式そのもの（$\CFORM$）」を個体（インスタンス）として扱う意味の場である\textbf{アスペクト的意味の場（Aspectual Field of Sense）}を導入し、任意の実質的意味の場の個別内容に一切立ち入ることなく、形式構造を基にした境界条件を純粋に記述し制限直積（ファイバー積）として結合できることを示す。第三に、変換元意味の場の形式とアスペクト的意味の場が与える境界条件の引き戻し（$\MSrc \times_{\CFORM} \MAspect$）から、変換先意味の場のスキーマを充足する状態（インスタンスモデル）を構成する一方向かつ不可逆な構造変換関手（\textbf{Transformation}）が、境界条件と変換先スキーマが与えられれば同型を除いて必ず一意に存在すること（\textbf{境界条件に基づく一意変換定理}）を証明するとともに、この境界値変換が数理物理学における\textbf{グリーン関数（Green's Function）による伝播作用素と完全に同型}であることを示す。第四に、「単一の意味の場内部の意味的関係（意味を伴う存在）」と「並立する意味の場同士を結ぶ背後のファンクター関係（現実世界での実在）」という\textbf{存在の二重構造}を峻別し、この形式主義的転回が自然科学と人文・社会科学の統一的基盤を与え、カントの超越論的哲学および西田幾多郎の「場所の論理・思考の普遍性」を厳密な数理形式として完成させることを論じる。
\end{abstract}

\section{序論（Introduction）}

\subsection{背景と問題設定：単一の意味の場の閉じたダイナミクスから、意味の場間の同期と次元転換へ}
前三稿~\cite{ota2026paper0,ota2026paper1,ota2026paper2} において、著者は関係を実在の基底とする存在論と複数の意味の場からなるユニバース $\mathcal{U}$ を定式化し、単一の意味の場における静的な意味構造や、普遍的実行セマンティクス $\Eexec$ のもとで物理学の理論記述と同型な動的体系を証明した。

しかし、現実の自然現象や工学的営為は孤立した単一の意味の場では完結せず、航空機の自動操縦や経済活動のように複数の意味の場が並立して強制的な同期や因果伝播を生じさせている。また、音声言語やプログラムコードへの変換といった、ある意味の場のスキーマを別の意味の場の制約に従うインスタンスへと構成する「次元転換作用素」の必然性と一意性を保証する体系が求められている。

\subsection{構造変換と意味の場間対応における4つの根源的問い}
本稿（Paper 3）では、以下の4つの問いに肯定的な数学的解答を与える。

\begin{enumerate}
    \item \textbf{ブリッジの自己閉包性と自己相似性の問い：} 並立する意味の場を結ぶ「ブリッジ」自体が独立した意味の場であり、Paper 1・2~\cite{ota2026paper1,ota2026paper2} の体系によって記述可能か。
    \item \textbf{形式への言及と「アスペクト的意味の場」の問い：} 概念モデリングの形式そのもの（$\CFORM$）を個体として扱うアスペクト的意味の場を導入し、実質的意味の場の個別内容に立ち入らず境界条件を記述・結合できるか。
    \item \textbf{境界条件に基づく一意の構造変換とグリーン関数同型の問い：} 引き戻しのもとで変換先スキーマを充足するインスタンスを構成する一方向かつ不可逆な構造変換関手が必ず一意に存在し、グリーン関数と同型で Shlaer--Mellor 法~\cite{shlaer1992,mellor2002} を厳密化できるか。
    \item \textbf{存在の二重構造と人文・自然科学の統一的基盤の問い：} 単一の場内部の意味的関係と並立する場を結ぶファンクター関係を峻別し、自然科学と人文・社会科学、カント~\cite{kant1781} や西田幾多郎~\cite{nishida1911,nishida1927} の哲学の統一的基盤を完成できるか。
\end{enumerate}

\subsection{本稿の構成}
本稿は以下の全7章で構成される。

\begin{itemize}
    \item \textbf{第2章（意味の場間の対応付けと「意味の場としてのブリッジ」）：} ブリッジが固有のスキーマ圏等を持つ独立した「意味の場」であることの定式化、プロジェクションクラス・ファンクター・対応関連クラスの導入、スキーマ制約による整合性判定と崩壊、強制同期およびブリッジ自身の同型定理と漸近的接近定理の証明。
    \item \textbf{第3章（形式への言及と「アスペクト的意味の場」）：} 実質的意味の場圏 $\CSUB$ と形式圏 $\CFORM$ の厳密な分離、「形式そのもの」を扱うアスペクト的意味の場およびアスペクト結合ブリッジ（制限直積）の確立。
    \item \textbf{第4章（境界条件に基づく一意の構造変換：Transformation）：} 引き戻し結合からの不可逆な構造変換関手 $\mathcal{T}$ の一意存在証明（境界条件に基づく一意変換定理および逆変換不可能性）、メディア選択の論理、第4.4節における Shlaer--Mellor 法のリカーシブデザインやモデルコンパイラ（Archetype 言語、T4 Template、MC-3020、C\#/Azure コンパイラ等）の数学的厳密化と限界の議論。
    \item \textbf{第5章（変換領域の諸定理とグリーン関数同型定理）：} 構造変換の4大構造定理の証明と、数理物理学のグリーン関数（Green's Function）による点源応答核による解の再構成と構造変換作用素の同型性の証明。
    \item \textbf{第6章（存在の二重構造と人文・自然科学の統一的基盤）：} 意味的関係と現実での実在の峻別、自然科学と人文科学の対称性、カント哲学の形式論的徹底と西田幾多郎の「場所の論理・思考の普遍性」の圏論的基礎づけの提示。
    \item \textbf{第7章（実践的意義と結論：Paper 4 への展望）：} 問題領域非依存のモデルコンパイラ、非機能要件の直交分離、生成系 AI 検証への工学的帰結の総括、最終稿（Paper 4）への展望。
\end{itemize}
\section{意味の場間の対応付けと「意味の場としてのブリッジ」}

\subsection{並立する意味の場と「意味的フレーズを持たない構造的対応」の必然性}
Paper 0~\cite{ota2026paper0} において、ユニバース $\mathcal{U}$ は単一の特権的形式空間によって上から統括されるのではなく、固有の静的スキーマ圏 $\mathcal{C}_i$（Paper 1~\cite{ota2026paper1}）と動的振る舞い（Paper 2~\cite{ota2026paper2}）を備えた複数の「意味の場（Fields of Sense）」$\mathcal{M}_1, \mathcal{M}_2, \dots$ が対等に並立する構造として定式化された。

ここで、並立する二つの意味の場 $\mathcal{M}_1$（スキーマ圏 $\mathcal{C}_1$）と $\mathcal{M}_2$（スキーマ圏 $\mathcal{C}_2$）が相互に連動する仕組み（ブリッジ）を厳密に記述しようとする際、単一の意味の場の内部で用いていた「意味的リレーションシップ（意味的フレーズ＋4種多重度）」をそのまま異なる意味の場のあいだに架けることは原理的に不可能である。なぜなら、Paper 0（構造崩壊定理）で証明した通り、日常言語やドメイン固有の「意味的フレーズ」は単一の意味の場の内部で意味を生成する一方で、異なる意味の場をまたぐと\textbf{言語的多義性（Polysemy）による構造崩壊（自明な恒等射への退化）}を引き起こすからである。

したがって、異なる意味の場同士を結ぶ対応付けにおいては、自然言語の述語（意味的フレーズ）を完全に排し、\textbf{「識別インデックスと両端の多重度のみによって定義される純粋な構造的対応射」}を採用しなければならない。

\subsection{ブリッジの自己相似性：「意味の場としてのブリッジ」と拡張道具立て}
本稿が確立する第一の核心命題は、\textbf{「二つの意味の場 $\mathcal{M}_1, \mathcal{M}_2$ を結ぶブリッジ $\mathcal{B}$ は、概念モデリング体系の外部から特別に持ち込まれる特権的な外部メカニズムではなく、それ自体が固有の概念情報モデル・状態モデル・アクションモデル・普遍的実行セマンティクスを備えた独立した第3の『意味の場』として構築できる」}という自己閉包性（自己相似性）である。

\begin{definition}[ブリッジの概念情報モデルにおける3つの拡張道具立て]\label{def:bridge-primitives}
二つの並立する意味の場 $\mathcal{M}_1 = (\mathcal{C}_1, \dots)$ と $\mathcal{M}_2 = (\mathcal{C}_2, \dots)$ を結ぶブリッジ $\mathcal{B}$ の概念情報モデル $\CBridge$ は、Paper 1 の基本プリミティブ（概念クラス・7種データ型・識別子・関連クラス）に加え、以下の3つの境界拡張要素によって構成される。

\begin{enumerate}
    \item \textbf{プロジェクションクラス（Projection Class：$P_A, P_B$）：}
    対向する意味の場 $\mathcal{M}_1$ の概念クラス $A \in \mathcal{K}_1$（および $\mathcal{M}_2$ の概念クラス $B \in \mathcal{K}_2$）の\textbf{公開境界（識別子およびドメインファンクション・公開事象が露出する状態の内包）のみをブリッジの内部へ射影した代理クラス}である。プロジェクションクラス $P_A$ は対向ドメインの内部属性や非公開の関連構造を一切覗き見ることができず（Paper 0・Paper 2 注記~3.1 の非透過性）、ブリッジの内部からは「外部ドメインの公開境界に立ち現れる個体集合 $\mathcal{F}_1(A)$ の射影」として振る舞う。
    \item \textbf{ファンクター（Functor：$F_k : P_A \leftrightarrow P_B$）：}
    ブリッジ内部の二つのプロジェクションクラス $P_A$ と $P_B$ を結ぶ構造的対応線である。単一の意味の場内部のリレーションシップと明確に区別するため、これを圏論的関手にちなんで\textbf{ファンクター（Functor）}と呼称する。ファンクター $F_k$ は\textbf{「一意のインデックス番号（識別ラベル $F_k$）と両端の多重度（\texttt{1}, \texttt{0..1}, \texttt{1..*}, \texttt{0..*}）のみを持ち、意味的フレーズを一切持たない」}純粋な構造的対応関係として定義される。
    \item \textbf{対応関連クラス（Associative Correspondence Class：$R_{AB}$）：}
    ファンクター $F_k$ によって結ばれる $P_A$ と $P_B$ のペアを実体化する関連クラス（Paper 1 公理~C-2 のスパン構造 $\langle p_A, p_B \rangle : R_{AB} \hookrightarrow P_A \times P_B$）である。対応関連クラス $R_{AB}$ には、対応付けの重み・変換パラメータ・同期履歴などのブリッジ固有の特徴値や、同期状態（未同期・同期済・競合中など）を管理する\textbf{状態モデルおよびアクションモデル（Paper 2）}をそのまま随伴させることができる。
\end{enumerate}
\end{definition}

\begin{theorem}[ブリッジの体系内閉包性と実行セマンティクス適用可能性]\label{thm:bridge-closure}
上記のプロジェクションクラス・ファンクター・対応関連クラスを備えたブリッジの概念情報モデル $\CBridge$ は、Paper 1 の商 2-圏の公理系（有限完備性・有限余完備性・スパン単射性・多重度制約）をすべて充足し、その上の状態モデルおよびデータフロー型アクションモデルは、Paper 2 の\textbf{普遍的実行セマンティクス $\Eexec$（形式圏 $\CFORM$）によって特別な例外規則を一切追加することなくそのまま駆動・実行可能}である。
\end{theorem}
\subsection{恣意的な対応付けの許容・両側スキーマ保存則・および強制的同期（\texttt{=>}, \texttt{<=>}, \texttt{REQUIRE}）}
現実世界やシステム設計を記述する際、人間は二つの意味の場 $\mathcal{M}_1, \mathcal{M}_2$ のあいだに\textbf{どのような対応付け（ブリッジ）であっても恣意的に記述すること自体は許容される}（これは Paper 2 において、\texttt{Can't Happen} 遷移やゼロ除算を含むモデル記述が形式的に許容されるのとまったく同型である）。しかし、そのブリッジが実際に意味の場として稼働・評価される際には、両側のスキーマ制約とのあいだで厳格な整合性判定が働く。

\begin{axiom}[両側スキーマ保存則と不整合時の始対象 $0$ への崩壊]\label{ax:bilateral-conservation}
ブリッジ $\mathcal{B}$ が二つの意味の場 $\mathcal{M}_1, \mathcal{M}_2$ を結合した合成系 $\mathcal{M}_1 \bowtie_{\mathcal{B}} \mathcal{M}_2$ において、任意の正規稼働ステップにおける状態組 $(\mathcal{F}_1, \mathcal{F}_{\mathcal{B}}, \mathcal{F}_2)$ は、接続された双方の意味の場の静的・動的制約およびブリッジ自身の制約を同時に充足しなければならない：
\begin{equation}
\mathcal{F}_1 \in \Valid(\mathcal{C}_1) \quad \wedge \quad \mathcal{F}_{\mathcal{B}} \in \Valid(\CBridge) \quad \wedge \quad \mathcal{F}_2 \in \Valid(\mathcal{C}_2).
\end{equation}
もしブリッジの恣意的な対応付けやアクションが、どちらか一方の意味の場において多重度違反・値域超過・\texttt{Can't Happen} 遷移を引き起こす場合、その対応付けは論理的帰結として\textbf{トポス内部論理の始対象 $0$（空空間・構造破綻）へと即時崩壊}する。
\end{axiom}

さらに、現実世界で観測される「飛行機の運行制御とニュートン力学の強制同期」や「契約締結と会計仕訳の強制連動」のように、一方の意味の場の状態変化が他方の意味の場の状態変化を\textbf{強制的に引き起こす（強制的に成り立つ対応付け）}メカニズムは、ブリッジの CIM 多重度と以下の3種の高次ブリッジ演算子によって厳密に記述される。

\begin{definition}[高次ブリッジ演算子：\texttt{=>}, \texttt{<=>}, \texttt{REQUIRE}]\label{def:bridge-operators}
ブリッジ $\mathcal{B}$ の状態モデル・アクションモデルおよびファンクター仕様においては、通常のドメイン内記述に加え、並立する意味の場間の因果伝播と拘束を宣言する3つの高次ブリッジ演算子が定義される（圏論的にはプロ関手・随伴関手・引き戻し極限に対応する）。
\begin{enumerate}
    \item \textbf{一方向因果伝播演算子（\texttt{=>}：Causal Propagation）：}
    $\mathcal{M}_1$ 側の公開事象の発生または状態到達をトリガーとして、ファンクター $F_k$ で結ばれた $\mathcal{M}_2$ 側の公開事象またはドメインファンクションを因果的に駆動する射（$\mathcal{E}_1^{\mathrm{public}} \Rightarrow \mathcal{E}_2^{\mathrm{public}} / F_{\mathrm{dom}, 2}$）。
    \item \textbf{双方向強制同期演算子（\texttt{<=>}：Forced Bidirectional Synchronization）：}
    ファンクター $F_k : P_A \leftrightarrow P_B$ の両端多重度 \texttt{1} と結合し、$\mathcal{M}_1$ 側の個体・状態変化と $\mathcal{M}_2$ 側の個体・状態変化が\textbf{同一の原子因果ステップ（Run-to-Completion）内で双方向に強制一致すること}を要求する同型拘束。一方の更新が他方のスキーマ制約に反して実行できない場合、系全体が始対象 $0$ へ崩壊する（例：航空機の操縦入力と物理的機体姿勢の強制同期）。
    \item \textbf{境界不変条件演算子（\texttt{REQUIRE}：Cross-Field Invariant Constraint）：}
    ファンクター $F_k$ で対応づけられたペア $(a, b) \in \mathcal{F}_{\mathcal{B}}(R_{AB})$ のあいだで、双方のドメインファンクションから射影される状態の内包（属性値や状態位相）が満たすべき事前・事後・不変論理式を宣言する引き戻し制約。違反時は即座に始対象 $0$ へ崩壊する。
\end{enumerate}
\end{definition}

\subsection{実在する対応関係の背後性と、ブリッジ自身の同型定理・漸近的接近定理}
ここで極めて重要な哲学的・数理的問いが生じる。抽象的なスキーマ $\mathcal{C}_1$ と $\mathcal{C}_2$ だけを現実世界から切り離して並べた場合、両者を結ぶブリッジの作り方は無数に考えられ、一意には決まらない。しかし、\textbf{対象となる現実世界の事象 $\mathcal{SYS}$（たとえば「大気中を飛行する実際の航空機システム」や「現実の商取引制度」）を固定したとき、その現実世界においてはすでに背後に隠れた真の対応関係（真のブリッジ $\mathcal{B}^*$）が秩序として実在し、現象を生み出している}。

そして、定理~\ref{thm:bridge-closure} によって「ブリッジそれ自体もまた独立した一つの意味の場である」ことが証明された以上、Paper 1（定理~5.3・5.4）および Paper 2（定理~5.3・5.4）で確立された\textbf{同型定理}と\textbf{漸近的接近定理}が、そのままブリッジ自身に対しても成立する。

\begin{theorem}[ブリッジの同型定理とアブダクションによる漸近的接近定理]\label{thm:bridge-iso-asymp}
現実世界の対象システム $\mathcal{SYS}$ における二つの意味の場 $\mathcal{M}_1, \mathcal{M}_2$ の相互作用を過不足なく厳密に記述する任意の二つのブリッジモデル $\mathcal{B}_1, \mathcal{B}_2$ は、プロジェクションクラスの粒度や表現記法に関わらず、\textbf{圏同値（同型：$\mathcal{B}_1 \simeq \mathcal{B}_2$）}となる。さらに、有限の人間が背後に隠れた真のブリッジ $\mathcal{B}^*$ を最初から見通せなくとも、観測される両ドメインの同期現象や不整合（始対象 $0$ への崩壊）からの\textbf{アブダクション}を通じて、プロジェクションクラスの分割・ファンクター多重度の精密化・\texttt{REQUIRE} 制約の追加を反復する近似列 $\mathcal{B}^{(0)} \to \mathcal{B}^{(1)} \to \cdots$ を構成することで、帰納的極限において真のブリッジ $\mathcal{B}^*$ へと漸近的に収束する（$\varinjlim \mathcal{B}^{(n)} \simeq \mathcal{B}^*$）。
\end{theorem}
\section{形式への言及と「アスペクト的意味の場（Aspectual Fields of Sense）」}

\subsection{実質的意味の場圏 $\CSUB$ と形式圏 $\CFORM$ の厳密な分離}
第2章で定式化した一般的なブリッジは、並立する二つの実質的な意味の場（たとえば「飛行制御の意味の場」と「ニュートン力学の意味の場」、あるいは「商取引の意味の場」と「財務会計の意味の場」）の\textbf{個体（インスタンス）と状態遷移}同士を結びつける意味の場であった。

これに対して、自然法則における対称性・保存則・時空幾何学の適用や、言語・図式・計算基盤といった表現空間への写像・制約付与においては、個々の物理粒子や個別事象ではなく、\textbf{「概念モデリングの『形式そのもの』――すなわち、概念クラス、特徴値、識別子、多重度、リレーションシップ、状態遷移表（STT）、およびデータフロー型アクション（DAG）――そのものを個体として扱い、その振る舞いや構造に形式的な境界条件を課す仕組み」}が不可欠となる。

\begin{definition}[実質的意味の場圏 $\CSUB$ と形式圏 $\CFORM$ の分離]\label{def:sub-vs-form}
特定の領域や現象の実質的な意味構造（内容）を記述する意味の場 $\mathcal{M} = (\mathcal{C}, \{M_A\}_{A \in \mathcal{K}}, \mathcal{DF})$ を\textbf{実質的意味の場（Substantive Field of Sense）}と呼び、それらがなす圏を\textbf{実質的意味の場圏（$\CSUB$）}と表記する。一方、Paper 2（定義~4.1）で導入した通り、あらゆる実質的意味の場が共有する構造的枠組み（概念クラス・特徴値・識別子・多重度・状態・事象・遷移・基本プロセス部品・ワイヤーという形式そのもの）を対象と射とする基準圏を\textbf{形式圏（$\CFORM$）}と呼ぶ。実質的意味の場圏 $\CSUB$ に属する任意の意味の場 $\mathcal{M}$ のスキーマ構造は、形式圏 $\CFORM$ 上のインスタンス関手 $\Phi_{\mathcal{M}} \in \Valid(\CFORM)$ として一意に表現される（構造的リフレクション）。
\end{definition}

\subsection{アスペクト的意味の場（Aspectual Field of Sense）の定式化}
この形式圏 $\CFORM$ を基盤として、実質的な内容を扱う $\CSUB$ の意味の場とは扱う対象次元を異にする意味の場である\textbf{アスペクト的意味の場（Aspectual Field of Sense）}を定義する。

\begin{definition}[アスペクト的意味の場 $\MAspect$]\label{def:aspectual-field}
\textbf{アスペクト的意味の場 $\MAspect = (\CAspect, \{M_{\alpha}\}, \mathcal{DF}_{\mathrm{Aspect}}, \Eexec)$} とは、以下の3条件を満たす「形式の次元」を扱う意味の場である。
\begin{enumerate}
    \item \textbf{形式の個体化（プロジェクション対象としての形式）：} $\MAspect$ においては、実質的意味の場の「概念クラス」「特徴値」「リレーションシップ」「状態」「事象」「アクション部品」といった\textbf{形式要素（$\CFORM$ の要素）そのものが、プロジェクションクラスを通じて『個体（インスタンス）』として立ち現れる}。
    \item \textbf{形式構造に基づく境界条件の記述：} $\MAspect$ は、時間離散化ステップ刻み幅 $\Delta t$ や数値精度、保存則・対称性制約、あるいは空間配置・排他調停・記憶構造といった\textbf{形式構造に基づく境界条件（Boundary Conditions）}を、固有の概念情報モデル・状態モデル・アクションモデルとして記述する。
    \item \textbf{実質的意味内容からの完全独立性（Substantive-Content Independence）：} $\MAspect$ は実質的意味の場圏 $\CSUB$ に属する個々の場の「形式（$\CFORM$）」のみを参照・制約し、対象となる実質的意味の場が力学現象であるか、経済取引であるか、生命現象であるかという\textbf{個別領域の実質的な意味内容には一切立ち入らない}。
\end{enumerate}
\end{definition}

\subsection{アスペクト結合ブリッジと制限直積（ファイバー積）としての状態構成}
従来の記述体系では、時間刻み幅や排他調停・空間配置といった境界条件が、実質的な現象や対象の本質的構造の記述の中に混入し、モデルの純粋性と普遍性を損なっていた。

アスペクト的意味の場 $\MAspect$ を導入することで、実質的意味の場 $\MSub \in \CSUB$ はその領域の本質的な意味秩序のみを純粋に記述し、外部環境や観測・実行空間が課すあらゆる境界条件は $\MAspect$ の側に完全に分離・局在化される。しかも、アスペクト的意味の場のスキーマ $\CAspect$ は概念モデリングの形式（$\CFORM$）に言及するものであるが、実際にその\textbf{状態（インスタンスモデル $\mathcal{F}_{\mathrm{Aspect}} \in \Valid(\CAspect)$）}が構成される際には、具体的な実質的意味の場 $\MSub$ の構造および状態（$\mathcal{F}_{\mathrm{Sub}}$）と対応関係を持つ。

したがって論理的必然として、対象となる実質的意味の場 $\MSub$ とアスペクト的意味の場 $\MAspect$ のあいだには、両者を結びつける専用の\textbf{アスペクト結合ブリッジ $\BAspect$} が定義される。このブリッジの構造は第2章の一般ブリッジとまったく同一の体系に従い、次のように一意に規定される。

\begin{proposition}[アスペクト結合ブリッジのスキーマと制限直積としての状態]\label{prop:aspect-coupling}
実質的意味の場 $\MSub \in \CSUB$ とアスペクト的意味の場 $\MAspect$ を結ぶアスペクト結合ブリッジ $\BAspect$ において：
\begin{enumerate}
    \item \textbf{ブリッジのスキーマ：} 形式圏 $\CFORM$（対象となる実質的意味の場の形式構造）とアスペクト的意味の場のスキーマ $\CAspect$ を基に定義され、両者のプロジェクションクラスをファンクターで結合する関連クラスとして構成される。
    \item \textbf{ブリッジの状態（インスタンスモデル）：} 対象となる実質的意味の場の状態 $\mathcal{F}_{\mathrm{Sub}}$ とアスペクト的意味の場の状態 $\mathcal{F}_{\mathrm{Aspect}}$ の、形式圏 $\CFORM$ 上の共通対応条件で制限された直積（ファイバー積／引き戻し）として構成される：
    \begin{equation}
    \mathcal{F}_{\mathrm{Bridge}} \cong \mathcal{F}_{\mathrm{Sub}} \times_{\CFORM} \mathcal{F}_{\mathrm{Aspect}}.
    \end{equation}
    \item \textbf{直交結合律：} 複数の直交するアスペクト的意味の場 $\mathcal{M}_{\mathrm{Aspect}, 1}, \dots, \mathcal{M}_{\mathrm{Aspect}, k}$ を結合する場合、ファイバー積の結合律により順序に依存せず一意に定まる：
    \begin{equation}
    \mathcal{M}_{\mathrm{bound}} \equiv \MSub \times_{\CFORM} \mathcal{M}_{\mathrm{Aspect}, 1} \times_{\CFORM} \cdots \times_{\CFORM} \mathcal{M}_{\mathrm{Aspect}, k}.
    \end{equation}
\end{enumerate}
\end{proposition}

この結合により、$\MSub$ の内部スキーマを一切書き換えることなく、その形式要素と状態に対して境界条件が直交的に付与される。
\section{境界条件に基づく一意の構造変換（Transformation）}

\subsection{変換の本質：ある意味の場の「スキーマ（形式）」から別の意味の場の「状態（インスタンス）」への構成}
前章までにおいて、並立する意味の場を結ぶ「意味の場としてのブリッジ」（第2章）と、概念モデリングの形式そのもの（$\CFORM$）を個体として扱う「アスペクト的意味の場」（第3章）を定式化した。本章では、これらを統合し、本体系における\textbf{構造変換（Transformation）}の数学的・存在論的本質を定式化する。

一見すると、構造変換は単なる「モデルからコードへの自動変換（モデルコンパイラ）」のように見えるかもしれない。しかし、その数理的本質は特定の工学技術に限定されるものではない。\textbf{構造変換（Transformation）とは本来、「ある変換元の意味の場 $\MSrc$ のスキーマ構造（形式）を、形式圏 $\CFORM$ を通じて取り出し、境界条件を規定するアスペクト的意味の場 $\MAspect$ の拘束のもとで、別の変換先の意味の場 $\MTarg$ のスキーマ制約に従う『一つの正当な状態（インスタンスモデル $\mathcal{F}_{\mathrm{Targ}} \in \Valid(\CTarg)$）』として構成する次元転換作用素」}である。

\begin{definition}[構造変換の三つ組と次元転換ブリッジ]\label{def:transformation-triple}
構造変換系は、以下の3つの意味の場およびそれらを結ぶ変換ブリッジ $\BTrans$ によって構成される。
\begin{enumerate}
    \item \textbf{変換元意味の場（Source Field of Sense：$\MSrc = (\CSrc, \dots)$）：} 表現・変換の対象となる本質的な意味構造（スキーマ $\CSrc$ と動的規則）を保持する意味の場。
    \item \textbf{アスペクト的意味の場（Boundary-Condition Field：$\MAspect$）：} 変換にあたって課される境界条件（表現メディアの選択規則、非機能要件、物理的・環境的制約）を、形式圏 $\CFORM$ 上の制約として記述する意味の場。
    \item \textbf{変換先意味の場（Target Field of Sense：$\MTarg = (\CTarg, \dots)$）：} 変換後の表現や構造を受け止める固有の文法・空間構造（スキーマ $\CTarg$）を備えた意味の場。
\end{enumerate}

ここで変換ブリッジ $\BTrans$ は、引き戻し結合 $\MSrc \times_{\CFORM} \MAspect$ によって取り出された\textbf{変換元の「形式要素（クラス・特徴値・多重度・スパン・状態・DAG アクション）」をプロジェクション対象とし、それらを変換先意味の場 $\MTarg$ のスキーマ $\CTarg$ を充足する「インスタンス（状態）の構成要素」へと多重度 \texttt{1} のファンクター束によって一方向写像（\texttt{=>}）するブリッジ}である（なお、変換元に初期状態データ $\mathcal{F}_{\mathrm{Src}, 0}$ が付随する場合には、それも同時に変換先インスタンス状態における初期値として埋め込まれる）。
\end{definition}

\subsection{境界条件に基づく一意変換定理と表現メディア・コンパイラへの適用}
一般的な二つの実質的意味の場のあいだでは、現実の制約を与えない限り無数の非同型なブリッジが許容された（第2章）。これに対し、形式圏 $\CFORM$ 全体を網羅するアスペクト的意味の場 $\MAspect$（境界条件）と、変換先スキーマ $\CTarg$ への完全被覆（多重度 \texttt{1} のファンクター束）が与えられたとき、変換先に構成される状態（インスタンスモデル）は自由度をすべて失い、一意に確定する。

\begin{theorem}[境界条件に基づく一意変換定理：Uniqueness of Boundary-Conditioned Transformation]\label{thm:unique-transformation}
任意の健全な変換元意味の場 $\MSrc$（その形式表現を $\Phi_{\mathrm{Src}} \in \Valid(\CFORM)$ とする）に対し、完全な境界条件を規定するアスペクト的意味の場 $\MAspect$ および変換先意味の場のスキーマ $\CTarg$ が与えられたとき、引き戻し $\MSrc \times_{\CFORM} \MAspect$ から誘導される構造変換関手
\begin{equation}
\mathcal{T}_{\MAspect} : \Valid(\CFORM) \longrightarrow \Valid(\CTarg)
\end{equation}
および変換先意味の場に構成されるインスタンス状態 $\mathcal{F}_{\mathrm{Targ}}^* = \mathcal{T}_{\MAspect}(\Phi_{\mathrm{Src}}) \in \Valid(\CTarg)$ は、\textbf{同型（圏同値）を除いて一意に決定される}。
\end{theorem}

\begin{proof}
変換元 $\MSrc$ のスキーマと動的規則は、形式圏 $\CFORM$ 上の有限個の生成元（概念クラス・特徴値・識別子・スパン・状態・16種基本プロセス部品の DAG）の余極限として表現される。アスペクト的意味の場 $\MAspect$ は、$\CFORM$ のすべての生成元に対して多重度 \texttt{1} のファンクター対応（境界条件）を割り当てるため、左カン拡張（Left Kan Extension：$\mathrm{Lan}$）および引き戻しの普遍性により、生成元上の割り当ては変換先スキーマ $\CTarg$ の公理をすべて充足する単一のインスタンス状態 $\mathcal{F}_{\mathrm{Targ}}^* \in \Valid(\CTarg)$ へと一意に拡張される。もし境界条件に矛盾や欠落があれば、公理~\ref{ax:bilateral-conservation} により系は始対象 $0$ へ崩壊する。したがって、無矛盾な境界条件のもとでは変換先状態 $\mathcal{F}_{\mathrm{Targ}}^*$ が同型を除いて一意に立ち現れる。
\end{proof}

\begin{corollary}[表現メディアの選択とソフトウェア・モデルコンパイラの統一的導出]\label{cor:media-and-compilers}
定理~\ref{thm:unique-transformation} の構造変換は、人間のコミュニケーションにおける表現メディアの選択からソフトウェア工学のモデルコンパイラに至るまで、まったく同一の数理形式として適用される。
\begin{enumerate}
    \item \textbf{表現メディアへの変換（言語・テキスト・図式・音響・映像）：}
    ある概念モデル（$\MSrc$）を他者に伝達する際、「音声言語」「テキスト文章」「ダイアグラム図式」「動画・サウンド」はいずれも固有の文法・空間構造（スキーマ $\mathcal{C}_{\mathrm{Media}}$）を持つ変換先の意味の場である。伝達目的や認知制約という境界条件（アスペクト $\MAspect$）を与えることで、元のモデルの形式は「テキストという状態」や「図という状態」として一意に構成される。
    \item \textbf{ソフトウェア工学におけるモデルコンパイラ（形式 $\to$ 実行プラットフォーム状態）：}
    実行プラットフォーム（特定のプログラミング言語、設定ファイル群、RDB、分散ミドルウェア等）は、それ自体が厳密な文法・構造制約を持つ一つの意味の場（スキーマ $\mathcal{C}_{\mathrm{Platform}}$）である。モデル変換とは、アプリケーションの形式（$\CFORM$）をアスペクト的意味の場（非機能要件・境界条件）のもとで「プラットフォーム言語スキーマの正当なインスタンス状態（ソースコード・設定群）」へと構成する操作にほかならない（実行時の初期インスタンスが必要な場合は、元のインスタンスモデルがプログラム状態の初期値として埋め込まれる）。そして変換の起点は個別アプリの意味内容ではなく普遍的な「形式（$\CFORM$）」であるため、いったん構築された変換ルールは\textbf{任意の問題領域のアプリケーションに対してそのまま100\%再利用・適用可能}である。
\end{enumerate}
\end{corollary}

\subsection{多段階変換の合成法則と不変条件の保存定理}
表現や実装のプロセスにおいては、ある意味の場の形式をいったん中間表現の意味の場の状態へと変換し、さらにそれを別の意味の場の状態へと多段階に変換する連鎖が頻繁に行われる。

\begin{theorem}[構造変換の合成法則と不変条件保存]\label{thm:transformation-composition}
境界条件アスペクト $\mathcal{M}_{\mathrm{Aspect}, 1}$ による変換関手 $\mathcal{T}_1$（$\mathcal{M}_0$ の形式を $\mathcal{M}_1$ の状態へ構成する変換）と、境界条件アスペクト $\mathcal{M}_{\mathrm{Aspect}, 2}$ による変換関手 $\mathcal{T}_2$（$\mathcal{M}_1$ の形式を $\mathcal{M}_2$ の状態へ構成する変換）がともに健全（始対象 $0$ へ崩壊しない）であるとき、その合成変換 $\mathcal{T}_2 \circ \mathcal{T}_1$ は、結合された複合アスペクト $\mathcal{M}_{\mathrm{Aspect}, 1} \bowtie \mathcal{M}_{\mathrm{Aspect}, 2}$ が誘導する一意変換関手と自然同型となる。さらに、変換元の正規稼働軌道 $\mathcal{F}_n \to \mathcal{F}_{n+1}$ は、変換先の意味の場における状態発展においても対応する不変条件と因果順序を忠実に保存する（可換図式の成立）。
\end{theorem}

\begin{corollary}[構造変換の不可逆性と逆変換不可能性定理]\label{cor:irreversibility}
構造変換関手 $\mathcal{T}_{\MAspect}$ は一方向（\texttt{=>}）の次元転換作用素であり、一般に単射（可逆）ではない。すなわち、変換先意味の場に構成された状態 $\mathcal{F}_{\mathrm{Targ}}^* \in \Valid(\CTarg)$ およびターゲットスキーマ $\CTarg$ のみから、元の変換元スキーマ $\CSrc$ および初期インスタンス状態 $\mathcal{F}_{\mathrm{Src}, 0}$ を\textbf{一意に逆算・復元することは原理的に不可能}である（Paper 0 の認識論的限界定理との構造的同型~\cite{ota2026paper0}）。これは、変換プロセスにおいて境界条件アスペクトによる構造の縮約・忘却が生じるとともに、変換元の「スキーマ」と「初期状態」が変換先では同一平面上の単一のインスタンス状態へと平坦化されるためである（ソフトウェア工学において、生成された実装コードから元の本質的な概念モデルを一意にリバースエンジニアリングできないことの数学的根拠）。
\end{corollary}

\subsection{工学的応用：Shlaer--Mellor 法におけるリカーシブデザイン（Recursive Design）の数学的厳密化}
定理~\ref{thm:unique-transformation}（境界条件に基づく一意変換定理）および系~\ref{cor:irreversibility}（逆変換不可能性）の最も洗練された工学的実現が、Shlaer--Mellor 法（および Executable UML）における\textbf{リカーシブデザイン（Recursive Design）}と\textbf{モデルコンパイラ（Model Compiler）}である~\cite{shlaer1992,mellor2002}。本体系の圏論的形式により、そのメカニズムは経験則から厳密な数学的定理へと昇華される。

\begin{enumerate}
    \item \textbf{ドメイン分離と再帰的適用（リカーシブ構造）の定式化：}
    リカーシブデザインでは、システム全体を単一の巨大モデルとして描くのではなく、アプリケーションドメイン（$\mathcal{M}_{\mathrm{App}}$）、サービスドメイン群（$\mathcal{M}_{\mathrm{Service}}$：ロギング、通信、永続化管理など）、そしてソフトウェアアーキテクチャドメイン（$\mathcal{M}_{\mathrm{Arch}}$：タスクスケジューリング、メモリ管理、分散トランザクション構造など）という複数の独立した意味の場（ドメイン）へと分離し、意味的フレーズを持たないファンクターによるブリッジで接続する。ここで「リカーシブ（再帰的）」と呼ばれる本質は、アーキテクチャドメインやサービスドメイン自身もまた特別な外部コードではなく、\textbf{アプリケーションドメインとまったく同じ形式圏 $\CFORM$（概念情報モデル・状態モデル・16種基本プロセス部品のアクションモデル）によって自己相似的にモデリングされる（定理~\ref{thm:bridge-closure}）}点にある。

    \item \textbf{アーキテクチャドメインによる形式の個体化（ポピュレーション）とモデルコンパイラの実装形式：}
    モデルコンパイルのプロセスにおいて、アーキテクチャドメイン $\mathcal{M}_{\mathrm{Arch}}$ はまさに第3章で定式化した\textbf{アスペクト的意味の場}として機能する。すなわち、アプリケーションドメイン $\mathcal{M}_{\mathrm{App}}$ のスキーマ構造（概念クラス、特徴値、リレーションシップ、状態遷移表、New OAL アクション部品）が、アーキテクチャドメイン側の形式クラス（\texttt{Class}, \texttt{Attribute}, \texttt{State}, \texttt{Event} 等）の\textbf{インスタンス状態（ポピュレーション：Population）}として対応づけられ、命題~\ref{prop:aspect-coupling} のアスペクト結合ブリッジ（制限直積 $\mathcal{F}_{\mathrm{Bridge}} \cong \mathcal{F}_{\mathrm{App}} \times_{\CFORM} \mathcal{F}_{\mathrm{Arch}}$）が形成される。

    そして、この変換ブリッジの内部に記述された\textbf{変換アクションモデル}（ポピュレートされた形式インスタンス群を走査し、アーキテクチャ固有の境界条件に従ってターゲット言語の構文要素へ写像する規則群）が普遍的実行セマンティクス $\Eexec$ のもとで駆動することで、ターゲット言語スキーマ $\mathcal{C}_{\mathrm{Lang}}$ を充足するソースコードおよび設定ファイル群（単一の変換先状態 $\mathcal{F}_{\mathrm{Lang}}^*$）が一意に構成される。

    現実のエンジニアリングにおいては、この変換アクションモデルを記述・実行する具体的な変換ルール記述言語やテンプレートエンジンとして、BridgePoint xtUML における \textbf{Archetype（アーキタイプ）言語}~\cite{xtumlorg} や、.NET 環境における \textbf{T4 Template（Text Template Transformation Toolkit）} などが用いられる。また、この変換アクションモデルを特定の実行環境向けに体系化したものが\textbf{モデルコンパイラ}であり、たとえば組込み C/C++ 環境向けの標準モデルコンパイラである \textbf{MC-3020}~\cite{mellor2002,xtumlorg} をはじめ、著者自身が設計・構築した \textbf{C\# 向けモデルコンパイラ}や \textbf{Microsoft Azure（クラウド／IoT・Digital Twins 分散環境）向けモデルコンパイラ}~\cite{ota2023azure} などは、いずれも同一のアプリケーション形式（$\CFORM$）を異なるアーキテクチャアスペクト（境界条件）のもとで各ターゲット言語・クラウド基盤の状態へと写像する、グリーン関数伝播核の具体的な工学的実現にほかならない。

    \item \textbf{リバースエンジニアリングおよびラウンドトリップ工学の位置づけと限界：}
    従来の多くの開発ツールでは、生成されたコード（ターゲット言語スキーマ $\mathcal{C}_{\mathrm{Lang}}$ 上のインスタンス状態）を直接修正し、そこから元のモデルへと逆変換して同期させる「ラウンドトリップ工学（Round-Trip Engineering）」が試みられてきた。もちろん、実質的意味の場圏 $\CSUB$ の適切なスキーマ（要求スキーマや抽出用スキーマ）と対応するブリッジさえ定義すれば、生成されたコードや成果物から「要求に属すると思われる状態の内包」を取り出すこと自体は可能であり、また逆に、元のモデルから表現形式（View）のスキーマと境界条件を切り替えることで、用途ごとに適した多様な表現（ダイアグラム図式、構造化テキスト、仕様ドキュメント、実行コード等）を一方向に導出することも自由に行える。
    しかしながら、Paper 1（第5章）~\cite{ota2026paper1} で証明した静的スキーマの逆算不可能性に加え、本稿の系~\ref{cor:irreversibility}（構造変換の不可逆性）が示す通り、境界条件アスペクトによる構造の縮約・忘却と平坦化を経た後の成果物のみから、元の本質的なアプリケーションスキーマ $\mathcal{C}_{\mathrm{App}}$ とアーキテクチャアスペクトの分離構造を\textbf{一意に完全復元（逆構成）することは数学的に不可能}である。したがって、モデルの修正は常に上流の意味の場（アプリケーションドメインまたはアーキテクチャアスペクト）において行い、一方向（\texttt{=>}）の変換を再実行して各表現を導出することこそが、整合性を永遠に保つ唯一のプロセスである。
\end{enumerate}

\section{変換領域の諸定理とグリーン関数同型定理}

\subsection{構造変換（Transformation）の4大構造定理}
Paper 1（静的スキーマ圏 $\mathcal{C}$）および Paper 2（動的振る舞いと対称モノイド圏 $\mathcal{DF}$）と完全に対称に、アスペクト的意味の場と変換ブリッジに基づく構造変換の領域においても、以下の4大構造定理が成立する。

\begin{theorem}[変換表現的完備性定理：Transformational Expressive Completeness]\label{thm:trans-completeness}
任意の計算可能な構造保存モデル変換（および任意の実行プラットフォームへのコード生成・マッピング規則）は、形式圏 $\CFORM$ 上のアスペクト的意味の場 $\MAspect$（概念情報モデル・状態モデル・16種基本プロセス部品の DAG アクション）と変換ブリッジ $\BTrans$ のみによって過不足なく記述可能である。
\end{theorem}

\begin{theorem}[変換存在性と不変条件閉包定理]\label{thm:trans-existence}
無矛盾な境界条件を定める任意のアスペクト的意味の場 $\MAspect$ とターゲットスキーマ $\CTarg$ に対し、変換元意味の場の正規稼働軌道を実装ドメインの正規稼働軌道へと写す構造保存変換関手 $\mathcal{T}_{\MAspect}$ が構成的に存在し、変換後の実装モデルもまた対称モノイド圏と普遍的実行セマンティクス $\Eexec$ のもとで閉じた実行可能モデルとなる。
\end{theorem}

\begin{theorem}[変換同型定理：Transformational Isomorphism]\label{thm:trans-isomorphism}
同一の変換元意味の場 $\MSrc$ と同一の形式構造に基づく境界条件を過不足なく表現する二つのアスペクト的意味の場 $\mathcal{M}_{\mathrm{Aspect}, 1} \simeq \mathcal{M}_{\mathrm{Aspect}, 2}$ から誘導される二つの変換関手 $\mathcal{T}_1, \mathcal{T}_2 : \MSrc \to \MTarg$ は、変換ルールの記述順序や中間変数名の差異に関わらず、\textbf{自然同型（$\mathcal{T}_1 \cong \mathcal{T}_2$）}となる。
\end{theorem}

\begin{theorem}[変換の漸近的洗練定理：Asymptotic Refinement of Transformations]\label{thm:trans-asymptotic}
有限のエンジニアが最初から完璧な実行環境変換ルールを見通せなくとも、実運用で観測される性能ボトルネック・競合デッドロック・リソース枯渇（始対象 $0$ への崩壊境界）からの\textbf{アブダクション}によって、アスペクト的意味の場 $\MAspect^{(n)}$ に新たな境界条件アスペクト（キャッシュ分割、ロック細粒度化、非同期キュー化など）を段階的に引き戻し合成（$\times_{\CFORM}$）していくことで、変換モデル列は帰納的極限において理想的な最適実装変換 $\mathcal{T}^*$ へと漸近的に収束する。
\end{theorem}

\subsection{数理物理学におけるグリーン関数（Green's Function）との構造的同型定理}
本稿の数理的クライマックスとして、アスペクト的意味の場（境界条件）を介したモデル変換作用素 $\mathcal{T}_{\MAspect}$ が、数理物理学における\textbf{グリーン関数（Green's Function：境界値問題の伝播核）}と数学的に完全に同型であることを証明する。

数理物理学（電磁気学・量子力学・熱伝導方程式・弾性体力学など）において、ある空間領域 $\Omega$ 上の線形微分作用素 $L$ と境界 $\partial\Omega$ 上の境界条件 $B(u) = 0$ のもとで、任意の源泉項（外力・電荷密度・初期分布）$f(\xi)$ に対する場の応答 $u(x)$ を求める問題 $L u(x) = f(x)$ を考える。このとき、点 $\xi$ に置かれた単位点源（ディラックのデルタ関数 $\delta(x - \xi)$）に対する境界条件付きの応答関数 $G(x, \xi)$（すなわち $L_x G(x, \xi) = \delta(x - \xi)$ かつ $B(G) = 0$ を満たす\textbf{グリーン関数}）を求めれば、任意の源泉 $f(\xi)$ に対する解 $u(x)$ は、点源応答核 $G(x, \xi)$ と源泉 $f(\xi)$ の重ね合わせ（積分作用素）によって一意に再構成される：
\begin{equation}
u(x) = \int_{\Omega} G(x, \xi)\, f(\xi)\, d\xi.
\end{equation}

ここで最も重要な事実は、\textbf{「グリーン関数 $G(x, \xi)$ は、具体的な源泉分布 $f(\xi)$（個別の問題入力）には一切依存せず、空間の微分構造 $L$ と境界条件 $B$ のみによって一意に決まる普遍的な伝播核（Propagator / Integral Kernel）である」}という点である。

\begin{definition}[グリーン関数境界値変換の圏論的仕様 $\mathcal{G}_{\mathrm{Phys}}$]\label{def:green-spec}
数理物理学におけるグリーン関数による境界値解決の形式 $\mathcal{G}_{\mathrm{Phys}}$ は、以下の4つ組によって規定される。
\begin{enumerate}
    \item \textbf{源泉空間（Source Space：$\mathcal{S}_{\mathrm{src}}$）：} 任意の入力分布 $f$ を点源（デルタ基底 $\delta_\xi$）の線形結合・余極限として分解する表現空間。
    \item \textbf{境界条件多様体（Boundary Conditions：$\mathcal{B}_{\mathrm{bc}}$）：} 領域の境界 $\partial\Omega$ 上で解が満たすべき拘束条件（ディリクレ条件・ノイマン条件等）。
    \item \textbf{グリーン関数核（Green's Propagator Kernel：$G(x, \xi)$）：} 境界条件 $\mathcal{B}_{\mathrm{bc}}$ のもとで単位点源 $\delta_\xi$ をターゲット空間の応答へと写像する、源泉 $f$ に依存しない普遍的積分核（プロ関手 $G : \mathcal{S}_{\mathrm{src}}^{\mathrm{op}} \times \mathcal{T}_{\mathrm{targ}} \to \Set$）。
    \item \textbf{重ね合わせ再構成作用素（Superposition / Coend Operator：$\int^{\xi} G(x, \xi) \otimes f(\xi)$）：} 任意の源泉 $f$ に核 $G$ を作用させて一意の解 $u$ を生成する余エンド（Coend）積分作用素。
\end{enumerate}
\end{definition}

\begin{theorem}[グリーン関数同型定理：Isomorphism with Green's Function Propagators]\label{thm:green-isomorphism}
アスペクト的意味の場 $\MAspect$ を境界条件とし、変換元 $\MSrc$ の形式を変換先 $\MTarg$ のインスタンス状態へと構成する構造変換体系
\begin{equation}
\mathfrak{T}_{\mathrm{CM}} \equiv \left(\Phi_{\mathrm{Src}} \in \Valid(\CFORM),\; \MAspect,\; \BTrans \subset \CFORM^{\mathrm{op}} \times \CTarg,\; \mathcal{T}_{\MAspect}\right)
\end{equation}
は、数理物理学におけるグリーン関数境界値変換 $\mathcal{G}_{\mathrm{Phys}} = (\mathcal{S}_{\mathrm{src}}, \mathcal{B}_{\mathrm{bc}}, G(x, \xi), \int^\xi)$ と\textbf{圏論的に同値（構造的に完全に同型）}である（表~\ref{tab:green-correspondence}）。
\end{theorem}

\begin{table}[htbp]
\centering
\caption{構造変換体系 $\mathfrak{T}_{\mathrm{CM}}$ と数理物理学のグリーン関数形式 $\mathcal{G}_{\mathrm{Phys}}$ の同型対応}
\label{tab:green-correspondence}
\renewcommand{\arraystretch}{1.35}
\begin{tabularx}{\textwidth}{>{\raggedright\arraybackslash}p{0.31\textwidth} >{\raggedright\arraybackslash}p{0.29\textwidth} >{\raggedright\arraybackslash}X}
\toprule
\textbf{構造変換体系 $\mathfrak{T}_{\mathrm{CM}}$（本稿）} & \textbf{数理物理学のグリーン関数形式 $\mathcal{G}_{\mathrm{Phys}}$} & \textbf{数学的・圏論的同型対応} \\
\midrule
\textbf{1. 変換元の形式 $\Phi_{\mathrm{Src}}$}（形式圏 $\CFORM$ の生成元分解） & \textbf{源泉空間 $\mathcal{S}_{\mathrm{src}}$}（点源 $\delta_\xi$ の重ね合わせとしての源泉 $f(\xi)$） & 変換元のスキーマ構造を、形式圏のプリミティブ生成元（クラス・関係・状態・部品＝点源）の余極限として展開する \\
\textbf{2. アスペクト的意味の場 $\MAspect$}（境界条件アスペクト） & \textbf{境界条件 $\mathcal{B}_{\mathrm{bc}}$}（ディリクレ・ノイマン境界条件等） & 変換元の意味内容（$f$）には立ち入らず、変換に際して従うべき外部境界制約を規定する \\
\textbf{3. 変換ブリッジ $\BTrans$}（形式 $\to$ 変換先状態のファンクター核） & \textbf{グリーン関数核 $G(x, \xi)$}（境界条件付き点源応答プロパゲータ） & 源泉 $f$（変換元ドメイン）に依存せず、境界条件 $\MAspect$ のみから一意に定まる普遍的変換核（プロ関手） \\
\textbf{4. 構成された変換先状態 $\mathcal{F}_{\mathrm{Targ}}^* = \mathcal{T}_{\MAspect}(\Phi_{\mathrm{Src}})$} & \textbf{解の再構成作用素 $u(x) = \int G(x, \xi) f(\xi)\, d\xi$} & 変換核をすべての形式要素にわたって余極限・余エンド（Coend）合成し、変換先スキーマ上の正当な一意の「状態（インスタンスモデル）」を構成する \\
\bottomrule
\end{tabularx}
\end{table}

\begin{proof}
任意の変換元モデル $\MSrc$ は、形式圏 $\CFORM$ のプリミティブ生成元 $\xi \in \CFORM$（個々のクラス・特徴値・スパン・状態・基本プロセス部品）の余極限（形式的な点源展開 $\MSrc \cong \int^{\xi} \CFORM(-, \xi) \otimes \Phi_{\mathrm{Src}}(\xi)$）として表現される。アスペクト的意味の場 $\MAspect$（境界条件）と変換先スキーマ $\CTarg$ が与えられると、各プリミティブ形式要素 $\xi$ を変換先スキーマ側の対応構造 $x \in \CTarg$ へと写す変換ブリッジ $\BTrans$ は、プロ関手（双関手）$K_{\MAspect} : \CFORM^{\mathrm{op}} \times \CTarg \to \Set$ を定める。このとき、定理~\ref{thm:unique-transformation} の一意変換関手 $\mathcal{T}_{\MAspect}$ による変換先状態 $\MTarg^*$ の生成は、圏論における\textbf{余エンド（Coend）積分公式}
\begin{equation}
\MTarg^*(x) \cong \int^{\xi \in \CFORM} K_{\MAspect}(x, \xi) \times \Phi_{\mathrm{Src}}(\xi)
\end{equation}
によって厳密に与えられる。これは積分核（グリーン関数）$G(x, \xi)$ による線形作用素の解の再構成公式 $u(x) = \int G(x, \xi) f(\xi)\, d\xi$ の圏論的定義そのものであり（Mac~Lane~\cite{maclane1998}）、系~\ref{cor:media-and-compilers} で示した「変換ルールが個別アプリケーションの問題領域 $f(\xi)$ に依存せず再利用可能である」という性質は、グリーン関数 $G(x, \xi)$ の源泉独立性と完全に同一の数学的帰結である。よって $\mathfrak{T}_{\mathrm{CM}} \simeq \mathcal{G}_{\mathrm{Phys}}$ が成り立つ。
\end{proof}

\section{存在の二重構造と人文・自然科学の統一的基盤}

\subsection{「意味を伴う存在」と「現実世界での実在」の二重構造}
前章までにおいて、並立する意味の場同士を結ぶ「意味の場としてのブリッジ」（第2章）、形式そのものを扱う「アスペクト的意味の場」（第3章）、および境界条件に基づく一意の構造変換とグリーン関数同型性（第4章・第5章）を確立した。ここで本体系全体を振り返るとき、世界における「存在」と「関係」をめぐる極めて重要な哲学的・存在論的帰結が浮かび上がる。

Paper 0~\cite{ota2026paper0} および Paper 1~\cite{ota2026paper1} において、著者は単一の対象（モノ）ではなく「関係（スパン構造）」を実在の基底とする存在論を提唱した。しかし、Paper 0〜3 の定式化を統合することで、一見すると同じ「何かと何かの関係」に見えるものの背後に、実は\textbf{存在論的次元のまったく異なる二種類の「関係」}が厳密に区別されて存在していることが明らかになる。

\begin{enumerate}
    \item \textbf{第一の関係――単一の意味の場内部における「意味的リレーションシップ」（＝意味を伴う存在）：}
    ある一つの意味の場（スキーマ圏 $\mathcal{C}$）の内部に閉じて、意味的フレーズ（日常言語や専門領域の述語）・識別子・4種の多重度に従って無矛盾に成立している関係である。この内部スキーマの制約（$\mathcal{F} \in \Valid(\mathcal{C})$）を満たすことは、対象が\textbf{「意味を伴って存在すること（Meaningful Being）」}の必要十分条件である。たとえば、精巧に構築された純粋数学の公理系、閉じたシミュレーションモデル、あるいは一貫した物語世界は、単一の意味の場の中でこの第一の関係を完全に充足しているため、「意味を伴って存在」している。
    \item \textbf{第二の関係――並立する意味の場同士を結ぶ背後の「ファンクター的対応」（＝現実世界での実在）：}
    しかし、単一の意味の場の内部だけで完結している閉じた構造は、それだけでは「現実世界に実在している（Actual / Real Existence）」とは言えない。\textbf{この現実世界において「実在する」とは、物理的観測の意味の場、感覚経験の意味の場、日常言語の意味の場、法制度・経済の意味の場など、並立する複数の意味の場のあいだに背後に隠れたブリッジが張り巡らされ、意味的フレーズを持たない純粋な構造的対応（ファンクター）と強制同期（\texttt{<=>}, \texttt{REQUIRE}）のスキーマに従って、諸々の意味の場が同時に共鳴・現象していること}にほかならない。
\end{enumerate}

\begin{proposition}[存在の二重構造：Meaningful Being と Real Actuality の峻別]\label{prop:dual-being}
ユニバース $\mathcal{U}$ における任意の事象・対象の実在性は、次の二層の構造的充足によって規定される。
\begin{itemize}
    \item \textbf{意味的成立（Meaningfulness）：} 単一の意味の場のスキーマ圏 $\mathcal{C}_i$ における意味的リレーションシップの充足（$\mathcal{F}_i \in \Valid(\mathcal{C}_i)$）。
    \item \textbf{現実的実在（Actuality）：} 並立する複数の意味の場の背後に実在するブリッジ $\mathcal{B}$ のプロジェクションクラスとファンクター対応の充足（$\mathcal{F}_1 \bowtie_{\mathcal{B}} \mathcal{F}_2 \neq 0$）。
\end{itemize}
\end{proposition}

人間が日常の観測において直接目にしているのは、個々の意味の場が境界（ドメインファンクションや公開事象）に映し出す「現象（内包）」のみであり、複数の意味の場を背後で結びつけているブリッジのファンクター構造そのものは直接目に見えない。したがって、観測される現象からのアブダクションによって背後のブリッジの対応付けを発見・記述し、接続された双方のスキーマ制約と照合して検証することこそが、\textbf{「単なる閉じた物語（架空の意味）」と「現実世界における実在」を峻別し、実在の秩序を暴き出す営みそのもの}である。

\subsection{自然科学と人文・社会科学の対称性}
この「意味的リレーションシップ（第一の関係）」と「背後のブリッジによるファンクター関係（第二の関係）」の二重構造は、人類の知の二大領域である\textbf{自然科学}と\textbf{人文・社会科学}の根本的な違いと、同時に両者が同一の数理形式の上に統一される理由を鮮やかに照らし出す。

\begin{itemize}
    \item \textbf{自然科学（Natural Sciences）におけるブリッジの所与性と発見：}
    自然界においては、並立する意味の場（たとえば「運動学の意味の場」と「電磁気学・熱力学の意味の場」、あるいは「微視的量子状態の意味の場」と「巨視的観測機器の意味の場」）を結びつける背後のブリッジ（自然法則・物理定数・境界条件）は、人間の主観や合意とは無関係に\textbf{「すでに在る前提（所与の秩序）」}として実在し、現象を生み出している。したがって、自然科学の営みは、観測される現象（内包）からのアブダクションとシミュレーションによる演繹的検証を通じて、その「すでに背後に実在する真のブリッジ $\mathcal{B}^*$ と意味の場のスキーマ」を人間が一歩ずつ発見し、漸近的に近似していくプロセスである（定理~\ref{thm:bridge-iso-asymp}）。
    \item \textbf{人文・社会科学（Humanities \& Social Sciences）におけるブリッジの構築と共有：}
    これに対して、言語・法制度・貨幣経済・組織・文化といった人文・社会領域においては、単一の意味の場（たとえば「法文の論理体系」や「会計帳簿の閉じた体系」）だけでなく、それらを現実の物理的行動や社会的合意と結びつけるブリッジ（契約の執行、貨幣と財の交換、翻訳規則、制度的規範）そのものが、\textbf{人間自身が歴史的・間主観的に記述・構築し、共有することによって初めて立ち現れる実在}である。したがって、人文・社会科学においては、この「第一の関係（閉じた言説内部の意味）」と「第二の関係（他領域と現実を結びつけるブリッジの対応）」を厳密に区別しなければ、単なる言葉遊び（閉じた物語）と現実の社会制度の実在を混同し、人文の思想そのものが成立し得ない。
    \item \textbf{形式の同一性による統一：}
    このように、「すでに在るブリッジを発見する（自然科学）」か「ブリッジを構築し共有する（人文・社会科学）」かという生成の向きは対称的でありながら、そこで機能している\textbf{「意味の場のスキーマ」「背後のファンクター・ブリッジ」「境界条件アスペクト」という数理的形式そのもの（$\CFORM$）は両領域において100\%同一}である。
\end{itemize}

\subsection{カントの超越論的哲学の徹底と西田幾多郎「場所の論理・思考の普遍性」の数理的完成}
本体系（Papers~0--3）が一貫して貫いてきた核心的立場は、\textbf{「概念モデリングとは世界を記述するための基本的かつ普遍的な『形式（\emph{Forma}）』であり、世界の成り立ち・在り様そのものの形式であるがゆえに、具体的にどのような意味の場や対応関係が存在するかという個別の『内容（\emph{Matter}）』には一切立ち入らない」}という徹底した形式主義的転回である。この立場は、近代西洋哲学および近代日本哲学が到達した二つの頂点に対し、決定的な数理的完成をもたらす。

\begin{enumerate}
    \item \textbf{カント（Immanuel Kant）の超越論的哲学を超える一歩進んだ転回：}
    カントは『純粋理性批判』~\cite{kant1781} において、「人間は物自体（\emph{Ding an sich}）を直接認識することはできず、我々が認識できるのは『現象』のみである」と喝破し、現象を成立させる超越論的形式を探求した。しかし、カントはその形式の根拠を人間主観の「感性の純粋直観（時間・空間）」や「純粋悟性概念（カテゴリー表）」という心理的・人間中心的なブラックボックスに委ねてしまったため、なぜそれが客観的実在そのものと一致するのかという点で歯切れの悪さを残した。
    これに対し、本体系（Papers~0--3）は物理的時空間のアプリオリさえも前提から排除し（Paper 2 注記~1.1~\cite{ota2026paper2}）、カントが「悟性」と呼んだ曖昧な能力を、人間主観から切り離された透明な圏論的形式（商 2-圏・対称モノイド圏・形式圏 $\CFORM$・グリーン関数変換核）へと完全に解体した。すなわち、\textbf{「物自体はもとより、現象そのものを直接扱う手前において、我々が明示的かつ無矛盾に記述できるのは、現象が立ち現れるための成立条件としての『形式』のみである」}という、カントからさらに一歩進んだ明快な立場を確立した。

    \item \textbf{西田幾多郎（Kitaro Nishida）の「場所の論理」と思考の普遍性の数理的基礎づけ：}
    同時に、この形式は孤立した個人の頭脳（主観）の中に閉じこもった私的枠組みではない。西田幾多郎が『善の研究』の純粋経験から後期哲学の「場所の論理」において直観した通り~\cite{nishida1911,nishida1927}、まず個物（モノ）があって関係が生まれるのではなく、まず普遍的な\textbf{「場所（\emph{Basho} ＝ 意味の場・トポス）」}があり、その場所の自己限定（スキーマ制約）として個物や主客の対立が立ち現れる。
    西田が確信した「個人を超えた思考の普遍性」の正体は、決して東洋的な神秘主義ではない。それはまさに、\textbf{「人間たちが同一の意味の場のスキーマ（形式 $\mathcal{C}$）と、そのスキーマに従った状態空間（$\Valid(\mathcal{C})$）を共有し、背後のブリッジによって互いに構造的に共鳴し合っていること」}そのものとして、本体系により透明かつ厳密な数理構造として基礎づけられる。
\end{enumerate}

\begin{proposition}[現象の正体としての構造変換：意識・思考の意味の場への Transformation]\label{prop:phenomena-transformation}
人間が経験する「現象（Paper 0 の内包~\cite{ota2026paper0}）」とは、諸々の意味の場の形式（$\CFORM$）が、知覚や認知の境界条件アスペクトのもとで「人間の意識・思考という意味の場」のインスタンス状態へとグリーン関数変換（Transformation）された結果である。
\end{proposition}

\section{実践的意義と結論（Paper 4 への展望）}

\subsection{ソフトウェア工学と生成系 AI 時代への工学的帰結}
本稿で確立した理論は、ソフトウェア工学およびモデル駆動開発（MDD）に以下の指針を与える。
\begin{enumerate}
    \item \textbf{問題領域非依存のモデルコンパイラの数学的保証：} 定理~\ref{thm:green-isomorphism}（グリーン関数同型定理）等により、変換ブリッジ（グリーン関数核 $G$）を一度構築すれば、アプリケーションの問題領域に関わらず変換規則を100\%再利用可能である。
    \item \textbf{非機能要件の完全な直交分離：} 境界条件はアスペクト的意味の場に局在化されファイバー積で合成されるため、要求モデルの純粋性と保守性が保たれる。
    \item \textbf{生成系 AI によるコード生成の厳密な検証基準：} グリーン関数同型定理により、AI によるコード生成を構造保存変換関手の近似計算として定式化し、妥当性を厳密に検証する基盤を提供する。
\end{enumerate}

\subsection{結論と最終稿（Paper 4）への架け橋}
本稿は単一の意味の場の理論を複数へと拡張し、ブリッジ、アスペクト的意味の場、グリーン関数同型、存在の二重構造という4つの柱を確立した。全四部作（Papers~0--3）により概念モデリングの数理的・工学的道具立てが完結し、最終稿 \textbf{Paper 4} ではこれらを統合してホログラフィー原理との同型性および存在論的必然性を証明する。

\begin{thebibliography}{99}
\bibitem{ota2026paper0} Ota, H. (2026). \emph{Relation-Based Ontology and a Category-Theoretic Formalization of Fields of Sense: Structural Collapse via Linguistic Polysemy and Phenomena as the Intension of a Universe} Zenodo. \url{https://doi.org/10.5281/zenodo.24267213}
\bibitem{ota2026paper1} Ota, H. (2026). \emph{Category-Theoretic Foundations of Conceptual Information Modeling: Rigorous Schema Representation, Isomorphism, and Asymptotic Refinement under Cognitive Finitude} Zenodo. \url{https://doi.org/10.5281/zenodo.24359112}
\bibitem{ota2026paper2} Ota, H. (2026). \emph{Category-Theoretic Foundations of Dynamic Modeling: State Transitions, Data-Flow Action Models, Universal Execution Semantics, and Isomorphism with Physical Field Theories} Zenodo. \url{https://doi.org/10.5281/zenodo.24430915}
\bibitem{maclane1998} Mac Lane, S. (1998). \emph{Categories for the Working Mathematician}. Springer.
\bibitem{kan1958} Kan, D. M. (1958). Adjoint Functors. \emph{Transactions of the AMS}, 87(2), 294--329.
\bibitem{benabou1973} B\'enabou, J. (1973). \emph{Les distributeurs}. Universit\'e Catholique de Louvain.
\bibitem{lawvere1963} Lawvere, F. W. (1963). Functorial Semantics of Algebraic Theories. \emph{PNAS}, 50(5), 869--872.
\bibitem{maclane1992} Mac Lane, S., \& Moerdijk, I. (1992). \emph{Sheaves in Geometry and Logic}. Springer.
\bibitem{joyal1991} Joyal, A., \& Street, R. (1991). The Geometry of Tensor Calculus, I. \emph{Advances in Mathematics}, 88(1), 55--112.
\bibitem{coecke2017} Coecke, B., \& Kissinger, A. (2017). \emph{Picturing Quantum Processes}. Cambridge University Press.
\bibitem{shlaer1992} Shlaer, S., \& Mellor, S. J. (1992). \emph{Object Lifecycles}. Yourdon Press.
\bibitem{mellor2002} Mellor, S. J., \& Balcer, M. J. (2002). \emph{Executable UML}. Addison-Wesley.
\bibitem{xtumlorg} xtUML.org. \emph{BridgePoint xtUML: Executable UML Modeling and Model Compilers}. \url{https://xtuml.org}
\bibitem{gabriel2015} Gabriel, M. (2015). \emph{Fields of Sense: A New Realist Ontology}. Edinburgh University Press.
\bibitem{kant1781} Kant, I. (1781/1787). \emph{Kritik der reinen Vernunft}. Cambridge University Press (1998).
\bibitem{nishida1911} Nishida, K. (1911/1990). \emph{Zen no Kenky\=u} (An Inquiry into the Good). Yale University Press.
\bibitem{nishida1927} Nishida, K. (1927). \emph{Hataraku Mono kara Miru Mono e} (Basho). Iwanami Shoten.
\bibitem{ota2023azure} Ota, H. (2023). \emph{Tutorial: Generated Library from Conceptual Model with Microsoft Azure Services}. GitHub. \url{https://github.com/kae-made/tutorial-generated-library-from-conceptual-model-with-azure-services}
\end{thebibliography}
\end{document}
